\section{Experimental design}
\label{sec:experimental-design}

\subsection{Hypotheses}
\label{sec:design-hypotheses}

The two statements below were written in the design document of the product before any wrapped run on a stress window, and are tested as written.

\begin{quote}
\textbf{H1.} HS-250, GARCH-t and FHS wrapped pass in at least one more stress window without worsening FZ0 by more than 5\%.\\
\textbf{H2 (refutation).} H1 is refuted if the wrapper fails independence in 2007--09.
\end{quote}

The terms are made exact by four conventions fixed with the hypotheses. A model \emph{passes} a window at a level when the Kupiec test and the conditional coverage test both have p-value at least 0.05; a test that the engine cannot evaluate (no hit in the window) passes nothing. The stress windows are the financial crisis window of 2007 to 2009, the pandemic window of 2020 and the rate-hiking window of 2022. H1 holds for a \emph{cell}, a model at a tail level (1\% or 2.5\%), when the wrapped model passes in at least one more of the three stress windows than the unwrapped one and its mean FZ0 is worse than that of the unwrapped model by at most 5\% in every stress window; it holds as a whole when it holds in all six cells and H2 does not occur. The wrapper \emph{fails independence} in a cell when the Christoffersen independence p-value of the wrapped model in the 2007--09 window is below 0.05. The step $\gamma = \gammaDefault$ and the band $[0.001, 0.25]$ were fixed before any run and were not tuned on any window.

Five further predictions follow from the theory and were written down before the generators of known truth were first run (they are listed in the header of the script \texttt{e5\_generators.py}): (a) in every scenario with a shift, for every model and level, the wrapped pass rate is above the unwrapped one; (b) with no shift the wrapped pass rate is within 0.05 of the unwrapped one and the median relative score change is within 2\%; (c) the wrapped mean hit rate is closer to the target in every scenario with a shift; (d) the repair is least complete in the shortest scenario, the decaying crash; (e) the wrapper lowers the mean absolute log error of the VaR against the true VaR for historical simulation in every scenario with a shift. The comparison with other recalibrations, the ablation and the extension from three models to ten were designed after the first results on the three models were known and carry no prediction.

\subsection{Data and its provenance}
\label{sec:design-data}

The data are three public daily series, taken as the case of Contraste takes them: the US equity market return from the Kenneth French data library (Mkt-RF plus RF), which is distributed for derived use only, so the returns are read from the device data root and only scores computed from them are committed; the ten-year US Treasury par yield from the Treasury's daily curve, converted to a return by the negative modified duration at the previous day's yield times the yield change; and the euro against the US dollar from the Federal Reserve's H.10 release. Six windows are built from them, with the portfolio of each window fixed in advance (\cref{tab:windows}): three stress windows (the first three rows), a foreign-exchange window of the financial crisis (251 days), a window of the bank turmoil of 2023 (64 days) and a calm control window of 2004 to 2006 (750 days). Equity returns are those of the market portfolio of the library; the mixed portfolios hold 50\% equity, 30\% rates and 20\% foreign exchange with daily rebalancing. A day on which any component is missing is skipped and counted, never filled. The hyperparameters of the models (the EWMA decay, the age-weighting decay, the threshold fraction of the extreme-value model) were selected once on the span 1995 to 2003, which ends before the earliest window, and are never changed. The factor library of \citet{fama1993} documents the construction of the equity series.

\begin{table*}[htbp]
\centering
\small
\caption{The six windows of case C11: the asset class of the portfolio, the first and last forecast day, the number of forecast days and of observed days of the common calendar, the public series of the portfolio and the role of the window in the evaluation.}
\label{tab:windows}
\begin{adjustbox}{max width=\textwidth}
\begin{tabular}{llllrrll}
\toprule
Window & Class & First forecast day & Last forecast day & Forecast days & Observed days & Series & Role \\
\midrule
Equity 2007--09 & equity & 2007-07-02 & 2009-06-30 & 504 & 1,504 & french--daily & stress \\
Equity 2020 & equity & 2020-02-03 & 2020-06-30 & 104 & 1,104 & french--daily & stress \\
Rates 2022 & rates & 2022-01-03 & 2022-12-30 & 249 & 1,249 & treasury--par & stress \\
FX 2008 & fx & 2008-07-01 & 2009-06-30 & 251 & 1,251 & h10--fx & other \\
Mixed 2023 & mixed & 2023-03-01 & 2023-05-31 & 64 & 1,064 & french--daily, treasury--par, h10--fx & other \\
Calm 2004--06 & mixed & 2004-01-02 & 2006-12-29 & 750 & 1,746 & french--daily, treasury--par, h10--fx & other \\
\bottomrule
\end{tabular}
\end{adjustbox}
\end{table*}


The three stress windows hold \hypStressDays{} forecast days in all (504, 104 and 249). The expected number of hits at the 1\% level is 5.0, 1.0 and 2.5 days, which is the reason the Kupiec test is read with its asymptotic caveat: its chi-square approximation is crude at counts this small, and the p-values of the engine are the asymptotic ones.

\subsection{Models compared, run by the product}
\label{sec:design-models}

The ladder holds eighteen models, all run by the pipeline of the product on the same windows with strictly rolling one-day forecasts that read only the returns before the forecast day: historical simulation over 250 and 500 days, the age-weighted variant, delta-normal and Cornish--Fisher with rolling moments, EWMA, GARCH, GJR-GARCH and EGARCH with Student-$t$ innovations, filtered historical simulation, the conditional extreme-value model, three CAViaR models (two for the VaR only, one with the ES step), three GAS-FZ models, and a quantile network trained on FZ0. Ten of them answer a VaR and an ES at any tail level and are wrapped: HS-250, HS-500, the age-weighted model, delta-normal, EWMA, GARCH-t, GJR-t, EGARCH-t, filtered historical simulation and the conditional extreme-value model. The Cornish--Fisher model and two CAViaR models give no ES, and the ES-CAViaR, the three GAS-FZ models and the network are fitted for one level and cannot be asked at another without a refit; they enter the ladder context and the score-space recalibrations but not the level wrapper.

Two recalibrations that act on the finished forecasts are run on the same forecasts as comparators. With the loss $L_t$ and the model's forecast $V_t$, define the score $S_t = L_t/V_t$, so that a hit is $S_t > 1$, and apply a scale $\theta_t$ to the VaR and the ES, $V_t^{\theta} = \theta_t V_t$ and $E_t^{\theta} = \theta_t E_t$.
The first is the rolling split-conformal scale: $\theta_t$ is the $k$-th smallest of the last $n \le \eqnRollWindow$ scores of the window with $k = \lceil (n+1)(1-\alpha)\rceil$, and 1 while $k > n$, clipped to $[0.25, 4]$ \citep{vovk2005,lei2018,barber2023}.
The second is the quantile tracker, the proportional term of conformal PID control \citep{angelopoulos2023}:
\begin{equation}
\theta_{t+1} = \max\bigl\{\theta_t + \eta\,(\indic\{S_t > \theta_t\} - \alpha),\ 0.25\bigr\}, \qquad \theta_1 = 1,
\label{eq:tracker}
\end{equation}
with $\eta = 0.05$ fixed before the run and $\eta \in \{0.01, 0.2\}$ reported as a sensitivity. Both are applied to all eighteen models; they need no model at another level.

\subsection{Metrics}
\label{sec:design-metrics}

For a sequence of $T$ forecast days with $x$ hits, the Kupiec test statistic is
\begin{equation}
\mathrm{LR}_{\mathrm{POF}} = -2\ln\frac{(1-\alpha)^{T-x}\alpha^{x}}{(1-x/T)^{T-x}(x/T)^{x}}, \label{eq:kupiec}
\end{equation}
asymptotically $\chi^2_1$ under the target frequency; the independence statistic $\mathrm{LR}_{\mathrm{ind}}$ is \cref{eq:lr-ind} computed with the observed transition counts of the sequence, asymptotically $\chi^2_1$, and the conditional coverage statistic is the sum $\mathrm{LR}_{\mathrm{POF}} + \mathrm{LR}_{\mathrm{ind}}$, asymptotically $\chi^2_2$ \citep{christoffersen1998,kupiec1995}. All three are computed by the engine and the p-values are the asymptotic ones. The score is the mean of \cref{eq:fz0} over the window, $\overline{\mathrm{FZ0}}$, and the change of a recalibration against the unwrapped model is
\begin{equation}
\Delta = \frac{\overline{\mathrm{FZ0}}_{\text{recalibrated}} - \overline{\mathrm{FZ0}}_{\text{unwrapped}}}{\bigl|\overline{\mathrm{FZ0}}_{\text{unwrapped}}\bigr|},
\label{eq:delta-fz0}
\end{equation}
positive when the recalibration scores worse, the quantity of the engine's \texttt{aci\_compare}. The ES is read through $Z_2$ of \cref{eq:z2}. The accuracy against a known truth, in the generators, is the mean absolute log error of the VaR against the true conditional VaR, $\frac{1}{T}\sum_t |\ln(V_t/V_t^{\mathrm{true}})|$, and its mean signed value as the bias. Four quantities describe the wrapper itself and are reported beside the tests: the hit rate $x/T$ with its Wilson interval, the mean working level, the number of days on which the band clipped the level, and the minimum of the path. The pass rate of a group of cells is the number of cells with a pass divided by the number of cells, so that it is a proportion over runs and carries a Wilson interval only where the runs are independent replications, as they are in the generators and are not in the real windows, where the cells share days.

\subsection{Protocol, statistics and planned ablations}
\label{sec:design-protocol}

On the real windows there is one realised path per window, so the uncertainty of a hit rate is its Wilson interval \citep{brown2001,wilson1927}, and the uncertainty of the change $\Delta$ is the 95\% percentile interval of a moving-block bootstrap with circular blocks of 10 days and 2,000 draws, which resamples the days of the paired daily scores jointly for the unwrapped and the recalibrated model \citep{kunsch1989,politis1994}. The seed of a draw is 20261104 plus a fixed integer per window, model and level. The bootstrap interval covers the variability of the score within the path and not the variability of the path itself, which one realised crisis cannot give; the generators supply the second kind. The generators of known truth use 100 replications per scenario, each with its own seed (the base 20261104 plus 1000 times the scenario index plus the replication), a design in the spirit of \citet{morris2019}: the pass rate of each method has its Wilson interval and the error against the truth its mean over replications.

The ablation changes one component of the wrapper at a time: the step ($\gamma \in \{0.001, 0.0025, 0.01, 0.02, 0.05\}$ against the default), the band (wider, $[10^{-5}, 0.5]$, and narrower, $[\alpha/2, 2\alpha]$) and the ES paired with the recalibrated VaR (the ES keeping the model's ratio to the VaR at the target level, and the ES left at the target level), on HS-250, GARCH-t and FHS at both levels in the three stress windows. The theory is checked by four numerical experiments (\cref{sec:reproducibility}): the identity of \cref{prop:identity} on the real level paths, the clustered sequence of \cref{prop:independence}, the scale sweep of \cref{prop:fixed-point} and the ES ratio of \cref{prop:es}. The unwrapped run of every model is checked against the forecasts of the ladder day by day, and the wrapped runs of the three models of H1 against the committed documents of the product; a mismatch stops the experiment.
